% Promoted after complete root and reduction/descent proof/code audits.
\subsection{Birational alternating pencils can have a reduced height-two base}
\label{subsec:abstract-Pfaffian-height-two-countermodel}

This subsection concerns abstract tensors over a field of characteristic
two.  It does not construct a Cartier tensor on an abelian surface.
In particular, the counterexample below does not assert failure of any
previously proved geometric monodromy statement.

\begin{proposition}[A uniform reduced Pfaffian counterexample]
\label{prop:abstract-birational-Pfaffian-height-two}
\status{Proved; independent audits passed.}
Let \(k\) have characteristic two, let \(m\ge2\), and put \(N=2m+1\).
There is an \(N\)-by-\(N\) alternating matrix \(B(x)\) of linear
forms in \(x_0,\ldots,x_{2m}\) with the following properties:
\begin{enumerate}
\item its generic rank is \(N-1\);
\item the exact common factor of its submaximal Pfaffians is
      \(x_0^{m-2}\);
\item after that factor is removed, its radical map is a birational
      quadratic transformation of \(\mathbf P^{N-1}\);
\item the reduced base support of that map is the union of a
      codimension-two linear space and a codimension-three linear space.
\end{enumerate}
Thus cancelling the full common factor and imposing birationality do
not imply that the reduced Pfaffian base has codimension at least three.
\end{proposition}
\begin{proof}
For \(1\le i\le2m\), let \(i^*\) be its partner in the consecutive
pairs \((1,2),(3,4),\ldots,(2m-1,2m)\).  Define \(B\) by
\[
 \begin{aligned}
 B_{0i}=B_{i0}&=x_i &&(1\le i\le2m),\\
 B_{2j-1,2j}=B_{2j,2j-1}&=x_0 &&(1\le j\le m-1),\\
 B_{2m-1,2m}=B_{2m,2m-1}&=x_1,
 \end{aligned}                                                   \tag{PH2.1}
\]
with all diagonal and other entries zero.  The lower \(2m\)-by-\(2m\)
block is invertible when \(x_0x_1\ne0\), so the alternating matrix
has generic rank \(2m=N-1\).

Let \(p_i\) be the Pfaffian obtained by deleting row and column \(i\).
If \(i>0\), the only nonzero perfect matching in that Pfaffian pairs
vertex zero with \(i^*\) and uses every remaining prescribed pair.
Consequently
\[
 \begin{aligned}
 p_0&=x_0^{m-1}x_1,\\
 p_i&=x_0^{m-2}x_1x_{i^*}
                  &&(1\le i\le2m-2),\\
 p_i&=x_0^{m-1}x_{i^*}
                  &&(i=2m-1,2m).
 \end{aligned}                                                   \tag{PH2.2}
\]
Divide by \(x_0^{m-2}\), and write the resulting vector as \(\sigma\):
\[
 \sigma_0=x_0x_1,\qquad
 \sigma_i=x_1x_{i^*}\ (1\le i\le2m-2),\qquad
 \sigma_i=x_0x_{i^*}\ (i=2m-1,2m).                               \tag{PH2.3}
\]
In particular \(\sigma_2=x_1^2\).  This coordinate and either last
coordinate have no common irreducible factor, so all the coordinates
are coprime.  The factor removed in (PH2.2) is therefore exact.

For independent coordinates \(y_0,\ldots,y_{2m}\), define
\[
 \tau_0=y_0^2,\qquad
 \tau_i=y_0y_{i^*}\ (1\le i\le2m-2),\qquad
 \tau_i=y_2y_{i^*}\ (i=2m-1,2m).                                \tag{PH2.4}
\]
Direct substitution gives polynomial identities
\[
       \tau(\sigma(x))=x_0x_1^2\,x,\qquad
       \sigma(\tau(y))=y_0^2y_2\,y.                              \tag{PH2.5}
\]
Both scalars are nonzero polynomials, proving birationality.

Finally, \(\sigma_2=0\) forces \(x_1=0\) set-theoretically.  The
only remaining equations are \(x_0x_{2m-1}=x_0x_{2m}=0\).  Thus
\[
 \operatorname{Supp}V(\sigma)=
 V(x_0,x_1)\ \cup\ V(x_1,x_{2m-1},x_{2m}).                       \tag{PH2.6}
\]
Both displayed linear spaces are nonempty and neither contains the
other.  Their codimensions are respectively two and three.  This
also computes the radical of the homogeneous base ideal, rather
than merely exhibiting a subset of its support.
\end{proof}

\begin{corollary}[The reciprocal contraction also attains its maximum]
\label{cor:abstract-Pfaffian-common-incidence}
\status{Proved; independent audits passed.}
The example arises from an alternating bilinear map
\(\beta:\Lambda^2k^N\to k^N\).  Its section contractions
\(M_s=\beta(s,-)\) have generic rank \(N-1\), and the actual
left-kernel map of \(M_s\) is birational, inverse to \(\sigma\).
There is an explicit common pair at which both \(M_s\) and
\(B_\lambda\) have rank \(N-1\).
\end{corollary}
\begin{proof}
Define \(\beta\) by
\(\lambda(\beta(s,t))=s^{\mathsf t}B(\lambda)t\), using the
coordinate pairing with the target.  This defines an alternating
map because every \(B(\lambda)\) is alternating.  Explicitly,
\[
 \begin{aligned}
 \beta_0(s,t)&=\sum_{j=1}^{m-1}
       (s_{2j-1}t_{2j}+s_{2j}t_{2j-1}),\\
 \beta_1(s,t)&=s_0t_1+s_1t_0+
       s_{2m-1}t_{2m}+s_{2m}t_{2m-1},\\
 \beta_i(s,t)&=s_0t_i+s_it_0\qquad(2\le i\le2m).
 \end{aligned}                                                   \tag{PH2.7}
\]
At \(s=e_0+e_2\), the output coordinates of \(M_s(t)\) are
\[
              (t_1,t_1,t_2+t_0,t_3,\ldots,t_{2m}),               \tag{PH2.8}
\]
so its rank is \(N-1\) and its unique left kernel is
\(\lambda=e_0^*+e_1^*\).  At this same \(\lambda\), every
prescribed lower pair in (PH2.1) has value one.  Therefore
\(B_\lambda\) also has rank \(N-1\), and its radical is \(ks\).

On the generic rank-\(N-1\) loci, the equation
\(M_s^{\mathsf t}\lambda=0\) is exactly \(B_\lambda s=0\).
The unique radical of \(B_\lambda\) is \(\sigma(\lambda)\).
The two birational maps (PH2.3)--(PH2.5) therefore identify
the left-kernel map of \(M_s\) with \(\tau\).
\end{proof}

\begin{remark}[Small matrix, exact scope, and certificate]
\label{rem:abstract-Pfaffian-height-two-scope}
\status{Proved; independent audits passed.}
For \(m=2\), no common factor is removed at all:
\[
 B(x)=
 \begin{pmatrix}
 0&x_1&x_2&x_3&x_4\\
 x_1&0&x_0&0&0\\
 x_2&x_0&0&0&0\\
 x_3&0&0&0&x_1\\
 x_4&0&0&x_1&0
 \end{pmatrix},\qquad
 \sigma=(x_0x_1,x_1x_2,x_1^2,x_0x_4,x_0x_3).                    \tag{PH2.9}
\]
For \(N=9\), the exact removed factor is \(x_0^2\), and the
same reduced height-two phenomenon persists.  Thus dimensions
divisible by nine do not themselves repair the abstract inference.

The actual geometric Cartier tensor has additional data not
established for this example: it is induced by the relative Cartier
pairing of a rank-four Frobenius pushforward bundle on an abelian
surface, has its theta-group equivariance, and has the specific
opposite-character Moore blocks and trivial-character tensor proved
earlier.  No identification with those data is asserted.
In particular this example does not prove reducibility of the actual
Lang incidence, failure of the full inertia theorem, or failure of
the proposed codimension bound for the actual Cartier tensor.

The retained standalone program
\path{code/monodromy/verify_birational_Pfaffian_height_two.py}
reconstructs the alternating matrices and their Pfaffians by the
perfect-matching recursion over \(\mathbf F_2\), independently checks
the exact gcd, both inverse compositions, the radical identities,
the two minimal monomial primes, and the explicit common-rank pair
for \(N=5,7,9,11,13\).  This finite exact supplement checks the
displayed formulas; the proof above, not extrapolation from those
five cases, proves the assertion for every odd \(N\ge5\).
The author run passed in \(0.003415901999687776\) seconds;
the unchanged, author-preserving independent root and reduction/descent
reruns passed in \(0.004231348997564055\) and
\(0.004291741002816707\) seconds, respectively.
Both reviewers independently rederived the complete uniform proof.
The root also checked the stated primary-source comparison; the
companion record makes no fresh original-source reading claim.
The exact candidate, all programs, outputs and audit records are retained.
\end{remark}

\paragraph{Primary-source boundary.}
Russo and Simis (2001), \emph{On Birational Maps and Jacobian Matrices},
Example 2.5, printed page 344 \cite[Example 2.5]{RussoSimis2001Birational}, assumes height three and its stated
local generator conditions before proving a Pfaffian map birational.
It does not prove the converse needed here.
Kustin, Polini, and Ulrich, \emph{The Equations Defining Blowup Algebras
of Height Three Gorenstein Ideals}, author version dated 14 October
2016, Theorem 9.1 \cite[Proposition 4.1 and Theorem 9.1]{KustinPoliniUlrich2016Blowups}, likewise assumes height three; Proposition 4.1 is
a birationality criterion, not an exclusion of height two.
Polishchuk (2022), \emph{Feigin--Odesskii brackets, syzygies, and
Cremona transformations}, Lemma 2.4.1 \cite[Theorem B and Lemma 2.4.1]{Polishchuk2022Cremona}, supplies a reciprocal
cofactor construction but no such base-codimension conclusion.
Its elliptic-curve application in Theorem B starts in characteristic
zero and already has the prescribed codimension-three secant locus.
The original statements and the indicated complete proofs were read;
none is imported as a theorem about the present Cartier tensor.
